\documentclass{article}
\usepackage[UTF8]{ctex}
\usepackage{amsmath,amssymb,amsfonts}
\usepackage{enumitem}
\usepackage{geometry}
\geometry{a4paper, margin=1in}
% === 错题本排版优化（不参与题目标识） ===
\setlength{\parindent}{0pt}
\setlength{\parskip}{0.35em plus 0.10em minus 0.08em}
\setlength{\emergencystretch}{3em}
\setlist{itemsep=0.25em, topsep=0.40em, parsep=0pt}
\allowdisplaybreaks[2]
% === 错题本排版优化结束 ===
\begin{document}
\section*{兰州大学2025年数学分析真题}
\section*{195．兰州大学2025年数学分析试卷}

\subsection{题目1}
1．\(\left\{x_{n}\right\}\) 的任意子列都有收敛到 0 的子列，则 \(\left\{x_{n}\right\}\) 是否收敛？为什么？

\subsection{题目2}
2．\(\sum_{n=1}^{\infty} a_{n}\) 条件收敛，则 \(\sum_{n=1}^{\infty} a_{n}^{3}\) 是否收敛？为什么？

\subsection{题目3}
1．求 \(\lim _{n \rightarrow \infty} \sum_{k=1}^{n} \frac{(2 k-1)^{4}}{n^{3}\left(n^{2}+k\right)}\) ．

\subsection{题目4}
2 ．求 \(\lim _{x \rightarrow 0}\left(\frac{\tan x}{x}\right)^{-x^{2}}\) ．

三．设 \(f\) 在 \([0,1]\) 上单调递增，\(f(0)>0, f(1)<1\) ．证明：存在 \(\xi \in(0,1)\) ，使得 \(f(\xi)=\xi^{2}\) ．

四．\(f(x), g(x)\) 在 \([0,+\infty)\) 上连续，\(\left|f^{\prime}(x)+f(x) g(x)\right| \leq|f(x)|\) ．证明：\(f(x) \equiv 0\) ．

五．\(f(x)\) 在 \((-\infty,+\infty)\) 上连续，\(f(x)\) 是周期函数，\(f\left(x^{2}\right)\) 也是周期函数，证明：\(f(x)\) 为常值函数．

六．\(f(x)\) 在［－1，1］上连续．

\begin{itemize}

  \item[1．] 证明： \(\int_{-1}^{1} f^{2}(x) \mathrm{d} x \geq \frac{1}{2}\left(\int_{-1}^{1} f(x) \mathrm{d} x\right)^{2}\) ．

  \item[2．] \(f(x)\) 单调递增，证明： \(\int_{-1}^{1} f^{2}(x) \mathrm{d} x \geq 2 f^{2}(0)\) 。
\end{itemize}

七．求 \(\frac{x^{2}}{3}+y^{2}+\frac{z^{2}}{2}=1\) 与 \(2 x+y+z=0\) 相交所得椭圆的面积．

八．\(\sum: z=x^{2}+y^{2}(0 \leq z \leq 1)\) ，方向取上侧，求 \(\iint_{\sum}(x+y) \mathrm{d} y \mathrm{~d} z+\left(y^{2}+z^{2}\right) \mathrm{d} z \mathrm{~d} x+\left(z^{3}+x^{3}\right) \mathrm{d} x \mathrm{~d} y\) ．

九．求 \(I(x)=\int_{0}^{\frac{\pi}{2}} \ln \left(1-x^{2} \cos ^{2} \theta\right) \mathrm{d} \theta, x \in(-1,1)\) ．

\subsection{题目12}
原题附加题：\(\left\{f_{n}(x)\right\}\) 在 \([0,1]\) 上连续，对任意的 \(x \in[0,1],\left\{f_{n}(x)\right\}\) 单调，\(\left\{f_{n}(x)\right\}\) 点态收敛于连续函数．证明：\(\left\{f_{n}(x)\right\}\) 一致收敛．
\end{document}
